% Eje sensible y sensibilidad transversal: una aceleración inclinada se reparte entre el eje sensible y el transversal.
\documentclass[tikz,border=4pt]{standalone}
\usepackage{resonancia}
\begin{document}
\begin{tikzpicture}[font=\sffamily\small, nota/.style={font=\sffamily\footnotesize, inner sep=1.5pt}]
  \draw[trazo, fill=rGrisClaro] (-0.7,-0.4) rectangle (0.7,0.4);
  \node[nota] at (-0.0,-0.65) {sensor};
  \draw[eje, color=rAzul, line width=1pt] (0,0.4) -- (0,3.2) node[above, nota, color=rAzul, align=center] {eje sensible\\(sensibilidad $S$)};
  \draw[eje, color=rVerde, line width=1pt, dash pattern=on 4pt off 2pt] (0.7,0) -- (3.4,0) node[right, nota, color=rVerde, align=left] {eje transversal\\(sensibilidad\\transversal $S_\mathrm{t}$)};
  \coordinate (A) at (2.3,2.0);
  \draw[fuerza] (0,0) -- (A) node[above right, nota, color=rNaranja] {aceleración $a$};
  \draw[punteada] (A) -- (0,2.0);
  \draw[punteada] (A) -- (2.3,0);
  \node[nota, color=rAzul, anchor=east] at (-0.05,1.25) {$a\cos\theta$};
  \node[nota, color=rVerde, anchor=north] at (1.5,-0.05) {$a\sin\theta$};
  \draw[trazo, color=rNaranja] (0,0.9) arc[start angle=90, end angle=41, radius=0.9];
  \node[nota, color=rNaranja] at (0.33,1.18) {$\theta$};
\end{tikzpicture}
\end{document}
